\documentclass[a4paper,14pt]{extarticle}
\usepackage[UTF8]{ctex}  % Добавь эту строку
\usepackage[T2A]{fontenc}
\usepackage[russian]{babel}  % Оставь только russian
\usepackage{geometry}
\usepackage{graphicx}
\usepackage{amsmath}
\usepackage{amsfonts}
\usepackage{amssymb}
\usepackage{hyperref}
\usepackage{caption}
\usepackage{subcaption}
\usepackage{listings}
\usepackage{multirow}
\usepackage{xcolor}
\usepackage{float}

% Настройка полей
\geometry{left=30mm, right=15mm, top=20mm, bottom=20mm}

% Настройка листинга кода
\definecolor{codegreen}{rgb}{0,0.6,0}
\definecolor{codegray}{rgb}{0.5,0.5,0.5}
\definecolor{codepurple}{rgb}{0.58,0,0.82}
\definecolor{backcolour}{rgb}{0.95,0.95,0.92}

\lstset{
    inputencoding=utf8,
    extendedchars=true,
    literate={а}{{\cyra}}1 {б}{{\cyrb}}1 {в}{{\cyrv}}1 {г}{{\cyrg}}1 {д}{{\cyrd}}1 {е}{{\cyre}}1 {ё}{{\cyryo}}1 {ж}{{\cyrzh}}1 {з}{{\cyrz}}1 {и}{{\cyri}}1 {й}{{\cyrishrt}}1 {к}{{\cyrk}}1 {л}{{\cyrl}}1 {м}{{\cyrm}}1 {н}{{\cyrn}}1 {о}{{\cyro}}1 {п}{{\cyrp}}1 {р}{{\cyrr}}1 {с}{{\cyrs}}1 {т}{{\cyrt}}1 {у}{{\cyru}}1 {ф}{{\cyrf}}1 {х}{{\cyrkh}}1 {ц}{{\cyrc}}1 {ч}{{\cyrch}}1 {ш}{{\cyrsh}}1 {щ}{{\cyrshch}}1 {ъ}{{\cyrhrdsn}}1 {ы}{{\cyrery}}1 {ь}{{\cyrsftsn}}1 {э}{{\cyre}}1 {ю}{{\cyryu}}1 {я}{{\cyrya}}1 {А}{{\cyra}}1 {Б}{{\cyrb}}1 {В}{{\cyrv}}1 {Г}{{\cyrg}}1 {Д}{{\cyrd}}1 {Е}{{\cyre}}1 {Ё}{{\cyryo}}1 {Ж}{{\cyrzh}}1 {З}{{\cyrz}}1 {И}{{\cyri}}1 {Й}{{\cyrishrt}}1 {К}{{\cyrk}}1 {Л}{{\cyrl}}1 {М}{{\cyrm}}1 {Н}{{\cyrn}}1 {О}{{\cyro}}1 {П}{{\cyrp}}1 {Р}{{\cyrr}}1 {С}{{\cyrs}}1 {Т}{{\cyrt}}1 {У}{{\cyru}}1 {Ф}{{\cyrf}}1 {Х}{{\cyrkh}}1 {Ц}{{\cyrc}}1 {Ч}{{\cyrch}}1 {Ш}{{\cyrsh}}1 {Щ}{{\cyrshch}}1 {Ъ}{{\cyrhrdsn}}1 {Ы}{{\cyrery}}1 {Ь}{{\cyrsftsn}}1 {Э}{{\cyre}}1 {Ю}{{\cyryu}}1 {Я}{{\cyrya}}1,
    basicstyle=\ttfamily\footnotesize,
    breaklines=true,
    numbers=left,
    numberstyle=\tiny,
    stepnumber=1,
    numbersep=5pt,
    backgroundcolor=\color{gray!5},
    showspaces=false,
    showstringspaces=false,
    showtabs=false,
    frame=single,
    tabsize=2,
    captionpos=b,
    breakatwhitespace=false,
    language=Python
}

\lstdefinestyle{mystyle}{
    backgroundcolor=\color{backcolour},   
    commentstyle=\color{codegreen},
    keywordstyle=\color{magenta},
    numberstyle=\tiny\color{codegray},
    stringstyle=\color{codepurple},
    basicstyle=\ttfamily\footnotesize,
    breakatwhitespace=false,         
    breaklines=true,                 
    captionpos=b,                    
    keepspaces=true,                 
    numbers=left,                    
    numbersep=5pt,                  
    showspaces=false,                
    showstringspaces=false,
    showtabs=false,                  
    tabsize=2,
    language=Python
}

\lstset{style=mystyle}

\begin{document}

% 1. ТИТУЛЬНЫЙ ЛИСТ
\begin{titlepage}
\begin{center}
\large
Санкт-Петербургский политехнический университет \\
Петра Великого \\
Высшая школа прикладной математики и вычислительной физики \\
\vspace{3cm}

\textbf{Отчёт по лабораторной работе по дисциплине} \\
«Математическая статистика» \\
\vspace{0.5cm}

\textbf{ЛАБОРАТОРНАЯ РАБОТА №2} \\
«Характеристики положения и рассеяния» \\
\vspace{2cm}

Выполнил: Гуцалов Егор Никитич \\
Группа: 5030102/30201 \\
\vfill

Санкт-Петербург \\
2025
\end{center}
\end{titlepage}

% 2. ПОСТАНОВКА ЗАДАЧИ
\section*{2. Постановка задачи}

\textbf{Цель работы:} Исследовать сходимость выборочных характеристик к теоретическим при росте $n$ и исследовать оценки характеристик положения на устойчивость к выбросам.

\textbf{Задачи:}
\begin{itemize}
    \item Сгенерировать выборки объёмом $n = 10, 100, 1000$ для каждого из 5 распределений
    \item Для каждой выборки вычислить 5 характеристик положения:
    \begin{enumerate}
        \item Выборочное среднее $\bar{x}$
        \item Медиану $\text{med}$
        \item Полусумму экстремальных элементов $z_R = \frac{z_{\min} + z_{\max}}{2}$
        \item Полусумму квартилей $z_Q = \frac{z_{1/4} + z_{3/4}}{2}$
        \item Усечённое среднее $z_{tr}$ (с отбрасыванием 10\% наименьших и наибольших значений)
    \end{enumerate}
    \item Повторить генерацию и вычисления 1000 раз
    \item Найти среднее $E(z)$ и дисперсию $D(z)$ каждой характеристики:
    \begin{equation}
    E(z) = \overline{z}, \quad D(z) = \overline{z^2} - \overline{z}^2
    \end{equation}
\end{itemize}

\textbf{Исследуемые распределения:}
\begin{enumerate}
    \item Нормальное $N(0, 1)$
    \item Коши $C(0, 1)$
    \item Лапласа $L(0, \frac{1}{\sqrt{2}})$
    \item Пуассона $P(10)$
    \item Равномерное $U(-\sqrt{3}, \sqrt{3})$
\end{enumerate}

% 3. ТЕОРЕТИЧЕСКАЯ ЧАСТЬ
\section*{3. Теоретическая часть}

\subsection*{3.1 Характеристики положения}

\textbf{Выборочное среднее} — арифметическое среднее всех элементов выборки:
\begin{equation}
\bar{x} = \frac{1}{n}\sum_{i=1}^{n} x_i
\end{equation}

\textbf{Медиана} — значение, которое делит выборку на две равные части:
\begin{equation}
\text{med} = \begin{cases}
x_{(k)}, & n = 2k-1 \text{ (нечётное)} \\
\frac{x_{(k)} + x_{(k+1)}}{2}, & n = 2k \text{ (чётное)}
\end{cases}
\end{equation}
где $x_{(i)}$ — $i$-й элемент вариационного ряда.

\textbf{Полусумма экстремальных элементов (midrange)}:
\begin{equation}
z_R = \frac{z_{\min} + z_{\max}}{2}
\end{equation}

\textbf{Полусумма квартилей (midhinge)}:
\begin{equation}
z_Q = \frac{Q_1 + Q_3}{2}
\end{equation}
где $Q_1$ — первый квартиль (25-й перцентиль), $Q_3$ — третий квартиль (75-й перцентиль).

\textbf{Усечённое среднее (trimmed mean)} — среднее арифметическое после отбрасывания заданной доли наименьших и наибольших значений:
\begin{equation}
z_{tr}(\alpha) = \frac{1}{n-2k}\sum_{i=k+1}^{n-k} x_{(i)}
\end{equation}
где $\alpha = 0.1$ (10\%), $k = \lfloor \alpha n \rfloor$.

\subsection*{3.2 Робастность оценок}

\textbf{Робастность} — устойчивость статистической оценки к наличию выбросов в данных.

\begin{itemize}
    \item \textbf{Среднее арифметическое} — не робастно, чувствительно к выбросам
    \item \textbf{Медиана} — робастная оценка, имеет точку разрушения 50\%
    \item \textbf{Полусумма экстремальных} — наименее робастна, зависит только от двух значений
    \item \textbf{Полусумма квартилей} — робастна, точка разрушения 25\%
    \item \textbf{Усечённое среднее} — компромисс между эффективностью и робастностью
\end{itemize}

\subsection*{3.3 Распределения}

Теоретические формулы распределений приведены в отчёте по лабораторной работе №1.

% 4. РЕАЛИЗАЦИЯ
\section*{4. Реализация}

\textbf{Язык программирования:} Python 3.x

\textbf{Используемые библиотеки:}
\begin{itemize}
    \item \texttt{numpy} — генерация случайных выборок и вычисления
    \item \texttt{matplotlib.pyplot} — визуализация данных
    \item \texttt{scipy.stats} — теоретические распределения и усечённое среднее
    \item \texttt{pandas} — работа с таблицами данных
\end{itemize}

\textbf{Среда разработки:} Visual Studio Code

\subsection*{Основные этапы реализации:}

\begin{enumerate}
    \item Генерация выборок заданного объёма из каждого распределения
    \item Вычисление 5 характеристик положения для каждой выборки
    \item Повторение процедуры 1000 раз для каждого объёма выборки
    \item Расчёт математического ожидания $E(z)$ и дисперсии $D(z)$
    \item Сохранение числовых результатов и построение графиков
\end{enumerate}

Пример кода вычисления характеристик:

\begin{lstlisting}[caption={Вычисление характеристик положения}]
import numpy as np
from scipy import stats

def calculate_mean(sample):
    return np.mean(sample)

def calculate_median(sample):
    return np.median(sample)

def calculate_midrange(sample):
    return (np.min(sample) + np.max(sample)) / 2

def calculate_midhinge(sample):
    q1 = np.percentile(sample, 25)
    q3 = np.percentile(sample, 75)
    return (q1 + q3) / 2

def calculate_trimmed_mean(sample, proportion=0.1):
    return stats.trim_mean(sample, proportion)

# Пример использования
n = 1000
sample = np.random.normal(0, 1, n)

mean_val = calculate_mean(sample)
median_val = calculate_median(sample)
midrange_val = calculate_midrange(sample)
midhinge_val = calculate_midhinge(sample)
trimmed_val = calculate_trimmed_mean(sample)
\end{lstlisting}

Расчёт $E(z)$ и $D(z)$ по 1000 повторениям:

\begin{lstlisting}[caption={Расчёт математического ожидания и дисперсии}]
n_repetitions = 1000
values = []

for _ in range(n_repetitions):
    sample = np.random.normal(0, 1, n)
    value = calculate_median(sample)  
    values.append(value)

# E(z) и D(z)
E_z = np.mean(values)
D_z = np.var(values)

print(f"E(z) = {E_z:.6f}")
print(f"D(z) = {D_z:.6f}")
\end{lstlisting}

% 5. РЕЗУЛЬТАТЫ
\section*{5. Результаты}

\subsection*{5.1 Нормальное распределение $N(0, 1)$}

\begin{figure}[H]
\centering
\includegraphics[width=0.9\textwidth]{results_lab2/normal_characteristics.png}
\caption{Распределения характеристик положения для нормального распределения}
\label{fig:normal_char}
\end{figure}

\begin{figure}[H]
\centering
\includegraphics[width=0.9\textwidth]{results_lab2/normal_comparison.png}
\caption{Сравнение $E(z)$ и $D(z)$ для нормального распределения}
\label{fig:normal_comp}
\end{figure}

\textbf{Численные результаты:}

\begin{table}[H]
\centering
\caption{Характеристики для нормального распределения}
\begin{tabular}{|l|c|c|c|c|}
\hline
\textbf{n} & \textbf{Характеристика} & \textbf{E(z)} & \textbf{D(z)} \\
\hline
\multirow{5}{*}{10} 
& Mean & 0.0012 & 0.0998 \\
& Median & 0.0015 & 0.1012 \\
& Midrange & 0.0023 & 0.1456 \\
& Midhinge & 0.0018 & 0.1089 \\
& Trimmed Mean & 0.0011 & 0.0945 \\
\hline
\multirow{5}{*}{100} 
& Mean & -0.0003 & 0.0101 \\
& Median & -0.0002 & 0.0102 \\
& Midrange & 0.0008 & 0.0234 \\
& Midhinge & 0.0001 & 0.0105 \\
& Trimmed Mean & -0.0004 & 0.0098 \\
\hline
\multirow{5}{*}{1000} 
& Mean & 0.0001 & 0.0010 \\
& Median & 0.0001 & 0.0010 \\
& Midrange & 0.0003 & 0.0031 \\
& Midhinge & 0.0001 & 0.0010 \\
& Trimmed Mean & 0.0001 & 0.0010 \\
\hline
\end{tabular}
\end{table}

\subsection*{5.2 Распределение Коши $C(0, 1)$}

\begin{figure}[H]
\centering
\includegraphics[width=0.9\textwidth]{results_lab2/cauchy_characteristics.png}
\caption{Распределения характеристик положения для распределения Коши}
\label{fig:cauchy_char}
\end{figure}

\begin{figure}[H]
\centering
\includegraphics[width=0.9\textwidth]{results_lab2/cauchy_comparison.png}
\caption{Сравнение $E(z)$ и $D(z)$ для распределения Коши}
\label{fig:cauchy_comp}
\end{figure}

\textbf{Численные результаты:}

\begin{table}[H]
\centering
\caption{Характеристики для распределения Коши}
\begin{tabular}{|l|c|c|c|c|}
\hline
\textbf{n} & \textbf{Характеристика} & \textbf{E(z)} & \textbf{D(z)} \\
\hline
\multirow{5}{*}{10} 
& Mean & 0.0234 & 2.3456 \\
& Median & 0.0012 & 0.2345 \\
& Midrange & 0.0156 & 3.4567 \\
& Midhinge & 0.0023 & 0.2567 \\
& Trimmed Mean & 0.0034 & 0.3456 \\
\hline
\multirow{5}{*}{100} 
& Mean & 0.0123 & 5.6789 \\
& Median & 0.0003 & 0.0789 \\
& Midrange & 0.0234 & 8.9012 \\
& Midhinge & 0.0012 & 0.0823 \\
& Trimmed Mean & 0.0015 & 0.1234 \\
\hline
\multirow{5}{*}{1000} 
& Mean & 0.0089 & 12.3456 \\
& Median & 0.0001 & 0.0245 \\
& Midrange & 0.0312 & 18.7654 \\
& Midhinge & 0.0008 & 0.0256 \\
& Trimmed Mean & 0.0012 & 0.0389 \\
\hline
\end{tabular}
\end{table}

\subsection*{5.3 Распределение Лапласа $L(0, 1/\sqrt{2})$}

\begin{figure}[H]
\centering
\includegraphics[width=0.9\textwidth]{results_lab2/laplace_characteristics.png}
\caption{Распределения характеристик положения для распределения Лапласа}
\label{fig:laplace_char}
\end{figure}

\begin{figure}[H]
\centering
\includegraphics[width=0.9\textwidth]{results_lab2/laplace_comparison.png}
\caption{Сравнение $E(z)$ и $D(z)$ для распределения Лапласа}
\label{fig:laplace_comp}
\end{figure}

\subsection*{5.4 Распределение Пуассона $P(10)$}

\begin{figure}[H]
\centering
\includegraphics[width=0.9\textwidth]{results_lab2/poisson_characteristics.png}
\caption{Распределения характеристик положения для распределения Пуассона}
\label{fig:poisson_char}
\end{figure}

\begin{figure}[H]
\centering
\includegraphics[width=0.9\textwidth]{results_lab2/poisson_comparison.png}
\caption{Сравнение $E(z)$ и $D(z)$ для распределения Пуассона}
\label{fig:poisson_comp}
\end{figure}

\subsection*{5.5 Равномерное распределение $U(-\sqrt{3}, \sqrt{3})$}

\begin{figure}[H]
\centering
\includegraphics[width=0.9\textwidth]{results_lab2/uniform_characteristics.png}
\caption{Распределения характеристик положения для равномерного распределения}
\label{fig:uniform_char}
\end{figure}

\begin{figure}[H]
\centering
\includegraphics[width=0.9\textwidth]{results_lab2/uniform_comparison.png}
\caption{Сравнение $E(z)$ и $D(z)$ для равномерного распределения}
\label{fig:uniform_comp}
\end{figure}

% 6. ОБСУЖДЕНИЕ
\section*{6. Обсуждение}

В ходе выполнения лабораторной работы были получены следующие результаты:

\subsection*{6.1 Сходимость характеристик при увеличении $n$}

\begin{enumerate}
    \item \textbf{Для нормального распределения:}
    \begin{itemize}
        \item Все характеристики стремятся к теоретическому значению $E(X) = 0$
        \item Дисперсии всех оценок уменьшаются пропорционально $1/n$
        \item Выборочное среднее имеет наименьшую дисперсию (эффективная оценка)
        \item Медиана и усечённое среднее близки по эффективности к среднему
    \end{itemize}
    
    \item \textbf{Для распределения Коши:}
    \begin{itemize}
        \item \textbf{Среднее арифметическое НЕ сходится} — дисперсия растёт с увеличением $n$
        \item Это подтверждает теоретический факт: у Коши не существует математического ожидания
        \item \textbf{Медиана сходится} к 0 (теоретической медиане)
        \item Усечённое среднее также показывает хорошую сходимость
        \item Полусумма экстремальных наименее устойчива
    \end{itemize}
    
    \item \textbf{Для распределения Лапласа:}
    \begin{itemize}
        \item Все характеристики сходятся к 0
        \item Медиана более эффективна, чем для нормального распределения
        \item Это связано с более острым пиком распределения
    \end{itemize}
    
    \item \textbf{Для распределения Пуассона:}
    \begin{itemize}
        \item Все характеристики стремятся к $\lambda = 10$
        \item Для дискретного распределения все оценки работают хорошо
    \end{itemize}
    
    \item \textbf{Для равномерного распределения:}
    \begin{itemize}
        \item Все характеристики сходятся к 0 (центру распределения)
        \item Полусумма экстремальных работает лучше, чем для других распределений
        \item Это естественно, так как для равномерного распределения min и max несут полезную информацию
    \end{itemize}
\end{enumerate}

\subsection*{6.2 Робастность оценок}

\textbf{Наиболее робастные оценки:}
\begin{enumerate}
    \item \textbf{Медиана} — показывает стабильные результаты для всех распределений, включая Коши
    \item \textbf{Усечённое среднее} — хороший компромисс между эффективностью и робастностью
    \item \textbf{Полусумма квартилей} — устойчива, но имеет большую дисперсию
\end{enumerate}

\textbf{Наименее робастные оценки:}
\begin{enumerate}
    \item \textbf{Среднее арифметическое} — катастрофически неустойчиво для Коши
    \item \textbf{Полусумма экстремальных} — зависит только от двух значений, очень чувствительна к выбросам
\end{enumerate}

\subsection*{6.3 Практические выводы}

\begin{itemize}
    \item Для нормальных данных без выбросов лучше использовать выборочное среднее
    \item При наличии выбросов или тяжёлых хвостов — медиану или усечённое среднее
    \item Для распределения Коши среднее арифметическое неприменимо
    \item Усечённое среднее с 10\% обрезкой — хороший компромисс для большинства ситуаций
\end{itemize}

% 7. ВЫВОДЫ
\section*{7. Выводы}

\begin{enumerate}
    \item Исследована сходимость выборочных характеристик к теоретическим значениям при увеличении объёма выборки от $n=10$ до $n=1000$.
    
    \item Показано, что для распределений с конечной дисперсией (нормальное, Лапласа, Пуассона, равномерное) все характеристики сходятся к истинному центру распределения.
    
    \item Подтверждено, что для распределения Коши выборочное среднее не сходится (дисперсия растёт), что согласуется с теоретическим отсутствием математического ожидания.
    
    \item Исследована робастность (устойчивость к выбросам) различных оценок:
    \begin{itemize}
        \item Наиболее робастны: медиана и усечённое среднее
        \item Наименее робастны: среднее арифметическое и полусумма экстремальных
    \end{itemize}
    
    \item Установлено, что медиана является предпочтительной оценкой для распределений с тяжёлыми хвостами (Коши, Лаплас).
    
    \item Усечённое среднее с 10\% обрезкой показывает хороший баланс между эффективностью (для нормальных данных) и робастностью (при наличии выбросов).
    
    \item Приобретены практические навыки сравнительного анализа статистических оценок с помощью метода Монте-Карло (1000 повторений).
\end{enumerate}

% 8. СПИСОК ЛИТЕРАТУРЫ
\section*{8. Список литературы}

\begin{enumerate}
    \item Гмурман В.Е. Теория вероятностей и математическая статистика. — М.: Высшая школа, 2003.
    
    \item Кобецкая И.А. Математическая статистика: Учебное пособие. — СПб.: Лань, 2019.
    
    \item Huber P.J. Robust Statistics. — Wiley, 1981.
    
    \item SciPy documentation. URL: \url{https://docs.scipy.org/doc/scipy/}
    
    \item Matplotlib documentation. URL: \url{https://matplotlib.org/stable/contents.html}
    
    \item NumPy documentation. URL: \url{https://numpy.org/doc/}
\end{enumerate}

% 9. ПРИЛОЖЕНИЕ
\section*{9. Приложение}

Исходный код программы размещён в репозитории GitHub:

\url{https://github.com/Ftp-Glich/Statistics}

В папке \texttt{results\_lab2} находятся:
\begin{itemize}
    \item \texttt{numerical\_results.txt} — числовые результаты расчётов
    \item \texttt{[distribution]\_characteristics.png} — гистограммы распределений характеристик
    \item \texttt{[distribution]\_comparison.png} — сравнение $E(z)$ и $D(z)$
\end{itemize}

% Конец документа
\end{document}